% Introduction

While synthetic data models like the Conditional Tabular GAN (CTGAN)~\cite{xu2019ctgan} offer a powerful workaround to strict data sharing frameworks such as GDPR~\cite{gdpr2016regulation}, HIPAA~\cite{hipaa1996act}, and the EU AI Act~\cite{euaiact2024}, proving the quality and safety of the generated data remains a fundamental challenge. Currently, a data officer runs local, isolated validation scripts and shares a self-certified quality report alongside the synthetic data. This fragile paradigm forces down-stream consumers to blindly trust the generating entity, lacking an immutable audit trail and frequently overlooking vital fairness audits~\cite{hardt2016equality}.

Distributed ledgers~\cite{nakamoto2008bitcoin} provide an avenue for tamper-proof provenance, but a blockchain cannot inherently verify the accuracy of the off-chain data being recorded---a limitation formally known as the ``oracle gap''~\cite{zheng2018blockchain}. This paper bridges that gap by cryptographically anchoring a rigorous, multi-metric objective evaluation system directly to the immutability of blockchain technology. 

To systematically explore this architecture, we structure our investigation around five central questions targeting: \textbf{Q1} (\textit{Efficacy}) objective trust-score improvements over single-party certification; \textbf{Q2} (\textit{Efficiency}) computational consensus overhead; \textbf{Q3} (\textit{Resilience}) Byzantine fault-tolerance limits; \textbf{Q4} (\textit{Scalability}) latency relative to dataset size; and \textbf{Q5} (\textit{Necessity}) the distinct isolation value of individual verification metrics.

Our primary contributions are:
\begin{enumerate}[leftmargin=*, itemsep=1pt]
  \item A \textbf{two-tier verification architecture} segregating objective machine-learning evaluation from tamper-proof cryptographic recording via Byzantine-tolerant median consensus.
  \item An \textbf{eleven-metric verification suite} spanning privacy (DCR, $k$-anonymity, MIA resilience, attribute disclosure), utility (Wasserstein, JS divergence, cross-correlation, TSTR ML efficacy), and fairness (four parity/impact constraints).
  \item A network-agnostic \textbf{hash-chain audit logger} paired with an \textbf{automated compliance engine} directly mapping quantitative scores to 15 specific mandates within GDPR, HIPAA, and the EU AI Act.
  \item An active \textbf{fairness post-processing module} that dynamically rectifies demographic imbalances during verification.
  \item A \textbf{comprehensive empirical evaluation} on the UCI Adult Income dataset~\cite{dua2019uci}, detailing multi-node consensus latency, sub-linear temporal scaling, and ablation testing.
\end{enumerate}
